\documentclass[10pt]{article}

% Specify the margins and text size.
\setlength{\textwidth}{7in}
\setlength{\textheight}{9.5in}
\setlength{\oddsidemargin}{0pt}
\setlength{\evensidemargin}{0pt}
\setlength{\topmargin}{0pt}
\setlength{\hoffset}{-.3in}
\setlength{\voffset}{-1in}

\setlength{\parskip}{5pt}
\setlength{\parindent}{0pt}

% Load some fonts and symbol packages
\usepackage{latexsym}
\usepackage{pifont}       % contains 'star' symbol for counterinsurgency handbook title
\usepackage{yfonts} 
\usepackage{amsmath}
\usepackage{amsfonts}

\usepackage{graphicx}     % actually, this is loaded by pstricks
\usepackage[T1]{fontenc}
\usepackage{ifthen}
\usepackage{pstricks,pst-grad,pst-text,pst-node,multido,pst-plot,calc,pst-3dplot}
%\usepackage[all]{xy}
%\usepackage{animate}

% The hyperref package inserts links into the pdf.
\definecolor{MyLinkColor}{rgb}{.1,.2,1}
\definecolor{MyCiteColor}{rgb}{.1,1,.2}
\definecolor{MyURLColor}{rgb}{.4,.4,.4}
\usepackage[backref=true,pagebackref=false,hyperindex,colorlinks=true,
  linkcolor=MyLinkColor,urlcolor=MyURLColor]{hyperref}


% The tweaklist package is something I found on the web.  It provides a simple interface
% for making changes to spacing used in the itemize and enumerate environments.  Comment
% this out if you don't care to use tweaklist.
\usepackage{tweaklist}
\renewcommand{\itemhook}{\setlength{\parskip}{2pt}\setlength{\parsep}%
{1pt}\setlength{\topsep}{0pt}\setlength{\itemsep}{0pt}}

\newcommand{\U}{\underline{\hspace{5pt}}}

\usepackage{listings}
\newcommand{\Z}{\hphantom{0}}

\begin{document}
\pagestyle{empty}
\lstset{language=R, showspaces=false, showstringspaces=false}

\href{http://www.su.edu}{\includegraphics[height=1.75cm]{sulogo.eps}}
\vspace{-1.69cm}

\hfill {\small
\begin{tabular}{cl}
& Math 207\\
& Introduction to Statistics\\
\end{tabular}
}
\setlength{\baselineskip}{1.05\baselineskip}
\bigskip

\begin{center}
\textbf{\large  Exam I}
\end{center}
No books, notes or calculators  are allowed on this exam.
\bigskip

1. (Hypothetical)  In a clinical trial, data collection usually starts at
\lq\lq baseline,\rq\rq\ when the subjects are 
recruited into the trial but before they are randomized to 
treatment or control.  Data collection continues until the end of the followup.
Two clinical trials on prevention of heart attacks report baseline data 
on weight, shown below.  In one of these trial, the randomization did not work.
Which  one and why? (15~points)\vspace{-5pt}
\newcommand{\Zc}{\hphantom{0,}}
\begin{center}
\begin{tabular}{cc}
\begin{tabular}{l}
\\
\\
\rput(-12pt,-3pt){(i)}
Treatment\vphantom{\Large Y}\\
Control\\[10pt]
\rput(-12pt,-3pt){(ii)}
Treatment\\
Control
\end{tabular}
&
\begin{tabular}{ccc}
Number of & Average &\\
persons & weight    & sd\\\hline
1,012   & 185 lb    & 28 lb\vphantom{\Large Y}\\
\Zc800 & 173 lb    & 26 lb\\[10pt]
1,006  & 202 lb    & 26 lb\\
1,002   & 162 lb    & 25 lb
\end{tabular}
\end{tabular}
\end{center}
\vspace{3.5in}

2. In a health study, subjects' heights has an average of 64 inches and sd of 3 inches.
Make a rough sketch of the histogram of subjects' heights.  About what percent of the subjects
were taller than 70 inches?
 (10~points)
\vfill
\eject

(20 points)

3.  a) Find the mean, median and sd of the list 141, 150, 148, 150, 150, 157, 154.
\vspace{2.6in}

{\ }\hphantom{n. }b) Which numbers in the list are within 1 SD of the mean?
\vspace{2.2in}


(15 points)

4. a) The Governor of California proposes to give all state employess a flat raise
of \$1,000 per year.  What would this do to the average annual salary of state employees?  
What would this do to the SD?
\vspace{2in}

{\ }\hphantom{n. }b) What would a 2.5\% increase in the salaries, across the board, do to the average
annual salary and to the SD?
\vfill
\eject


5. A histogram of monthly wages for part-time employees is shown below (densities are marked
in parentheses).  Nobody earned more than \$1,000 per month.  (20 points)

\begin{center}
\begin{pspicture}(0,-0.8)(10,2.3)
\psset{yunit=0.1}
\psframe(0,0)(1,10)\rput[b](0.5,10.1){(10)}
\psframe(1,0)(2,5)\rput[b](1.5,5.1){(5)}
\psframe(5,0)(10,5)\rput[b](7.5,5.1){(5)}
\psaxes[Dy=10](0,0)(10,20)
\rput{90}(-1.1,10){\% per \$100}
\rput(5,-10){Wages (hundreds of dollars)}
\end{pspicture}
\end{center}

\hspace{10pt} a) The block
over the class interval from \$200 to \$500 is missing. 
Determine how tall it must be then add it to the 
histogram.
\vspace{3in}

\hspace{10pt} b) About what percent of the employees earned between \$700 and \$900 per month?
\vspace{1.75in}

\hspace{10pt} c) About what percent of the employees earned more than \$200 per month?
\vfill
\eject



6. The table below shows results from two experiments designed to test the Salk vaccine in  1954.
(20 points)
\vspace{-5pt}
\begin{center}
{%\footnotesize
\begin{tabular}{cc}
Randomized controlled          & \\
double-blind experiment        & NFIP study\\
   \begin{tabular}{lcc}\hline
                 & Size      & Rate${}^*$\vphantom{\LARGE Y}\\
      \color{black}Treatment  & \color{black}200,000   & {\color{black}{28}}\\
      \color{black}Control    & \color{black}200,000   & {\color{black}{71}}\\
      No consent & 350,000   & 46\\
   \end{tabular}
&\hspace{20pt}
   \begin{tabular}{lcc}
\hline
                 & Size        &  Rate${}^*$\vphantom{\LARGE Y}\\
      \color{black}Grade 2 (vaccine)\hfill        & \color{black}225,000 & 25\\
      \color{black}Grades 1 and 3 (control)        & \color{black}725,000 & {\color{black}{54}}\\
      Grade 2 (no consent)\hfill     & 125,000 & {\color{black}{44}}\\
   \end{tabular}
\end{tabular}}\vspace{-4pt}
\end{center}
{${}^*$Rate is number of cases per 100,000 subjects.}
\smallskip

\hspace{10pt} a) Describe the treatment and control groups and how they were selected 
in the two experiments.  
\vspace{1.75in}


\hspace{10pt} b) What does \textit{consent} mean in the context of the experiments?
\vspace{1.25in}

\hspace{10pt} c) Which rates demonstrate the vaccine's effectiveness in the two experiments 
(there are two rates to compare for each experiment)?  
\vspace{1.25in}

\hspace{10pt} d) Which experiment gave stronger evidence for the vaccine's effectiveness and why?
\vspace{1.5in}

%\hspace{10pt} e) What does \textit{double-blind} mean?
%\vspace{1in}

\hspace{10pt} e) What impact did the different sizes of the six groups have on your conclusions and why?



\vfill
\eject
\end{document}

